Apply the following construction separately to \(E=A_Q\cap B_Q\) and \(E=B_Q\). Let \(\mathcal T_v\) be the finite set of signed-isotone complete local response tables specified by Definition \(\mathrm{def:signed\text{-}response\text{-}types}\), and let \(R_v\) denote the realized complete local response table. For every \(t=(t_v)_{v\in V}\in\prod_{v\in V}\mathcal T_v\), recursively evaluate the literals of \(E\) in topological order in \(\mathsf O_{G,M}(Q)\), after deleting every local restriction inconsistent with the normalized Boolean-lattice order. If \(e_E(t)\in\{0,1\}\) is the resulting event indicator, define
\[
 \mathsf N_{G,M}(Q;E)
 =\sum_{t:e_E(t)=1}\prod_{D\in\mathcal D(G)}
 \mathbf 1\{(R_v)_{v\in D}=t_D\},
\]
with terms ordered lexicographically by a fixed ordering of vertices and local tables. The two forms for \(A_Q\cap B_Q\) and \(B_Q\) constitute \(\mathsf N_{G,M}(Q)\). Thus the normal form is the expansion in the fixed exact response-type atom basis, rather than an unspecified choice among algebraically equivalent M\"obius expansions.